% Options for packages loaded elsewhere
\PassOptionsToPackage{unicode}{hyperref}
\PassOptionsToPackage{hyphens}{url}
\documentclass[
]{article}
\usepackage{xcolor}
\usepackage{amsmath,amssymb}
\setcounter{secnumdepth}{-\maxdimen} % remove section numbering
\usepackage{iftex}
\ifPDFTeX
  \usepackage[T1]{fontenc}
  \usepackage[utf8]{inputenc}
  \usepackage{textcomp} % provide euro and other symbols
\else % if luatex or xetex
  \usepackage{unicode-math} % this also loads fontspec
  \defaultfontfeatures{Scale=MatchLowercase}
  \defaultfontfeatures[\rmfamily]{Ligatures=TeX,Scale=1}
\fi
\usepackage{lmodern}
\ifPDFTeX\else
  % xetex/luatex font selection
\fi
% Use upquote if available, for straight quotes in verbatim environments
\IfFileExists{upquote.sty}{\usepackage{upquote}}{}
\IfFileExists{microtype.sty}{% use microtype if available
  \usepackage[]{microtype}
  \UseMicrotypeSet[protrusion]{basicmath} % disable protrusion for tt fonts
}{}
\makeatletter
\@ifundefined{KOMAClassName}{% if non-KOMA class
  \IfFileExists{parskip.sty}{%
    \usepackage{parskip}
  }{% else
    \setlength{\parindent}{0pt}
    \setlength{\parskip}{6pt plus 2pt minus 1pt}}
}{% if KOMA class
  \KOMAoptions{parskip=half}}
\makeatother
\usepackage{color}
\usepackage{fancyvrb}
\newcommand{\VerbBar}{|}
\newcommand{\VERB}{\Verb[commandchars=\\\{\}]}
\DefineVerbatimEnvironment{Highlighting}{Verbatim}{commandchars=\\\{\}}
% Add ',fontsize=\small' for more characters per line
\newenvironment{Shaded}{}{}
\newcommand{\AlertTok}[1]{\textcolor[rgb]{1.00,0.00,0.00}{\textbf{#1}}}
\newcommand{\AnnotationTok}[1]{\textcolor[rgb]{0.38,0.63,0.69}{\textbf{\textit{#1}}}}
\newcommand{\AttributeTok}[1]{\textcolor[rgb]{0.49,0.56,0.16}{#1}}
\newcommand{\BaseNTok}[1]{\textcolor[rgb]{0.25,0.63,0.44}{#1}}
\newcommand{\BuiltInTok}[1]{\textcolor[rgb]{0.00,0.50,0.00}{#1}}
\newcommand{\CharTok}[1]{\textcolor[rgb]{0.25,0.44,0.63}{#1}}
\newcommand{\CommentTok}[1]{\textcolor[rgb]{0.38,0.63,0.69}{\textit{#1}}}
\newcommand{\CommentVarTok}[1]{\textcolor[rgb]{0.38,0.63,0.69}{\textbf{\textit{#1}}}}
\newcommand{\ConstantTok}[1]{\textcolor[rgb]{0.53,0.00,0.00}{#1}}
\newcommand{\ControlFlowTok}[1]{\textcolor[rgb]{0.00,0.44,0.13}{\textbf{#1}}}
\newcommand{\DataTypeTok}[1]{\textcolor[rgb]{0.56,0.13,0.00}{#1}}
\newcommand{\DecValTok}[1]{\textcolor[rgb]{0.25,0.63,0.44}{#1}}
\newcommand{\DocumentationTok}[1]{\textcolor[rgb]{0.73,0.13,0.13}{\textit{#1}}}
\newcommand{\ErrorTok}[1]{\textcolor[rgb]{1.00,0.00,0.00}{\textbf{#1}}}
\newcommand{\ExtensionTok}[1]{#1}
\newcommand{\FloatTok}[1]{\textcolor[rgb]{0.25,0.63,0.44}{#1}}
\newcommand{\FunctionTok}[1]{\textcolor[rgb]{0.02,0.16,0.49}{#1}}
\newcommand{\ImportTok}[1]{\textcolor[rgb]{0.00,0.50,0.00}{\textbf{#1}}}
\newcommand{\InformationTok}[1]{\textcolor[rgb]{0.38,0.63,0.69}{\textbf{\textit{#1}}}}
\newcommand{\KeywordTok}[1]{\textcolor[rgb]{0.00,0.44,0.13}{\textbf{#1}}}
\newcommand{\NormalTok}[1]{#1}
\newcommand{\OperatorTok}[1]{\textcolor[rgb]{0.40,0.40,0.40}{#1}}
\newcommand{\OtherTok}[1]{\textcolor[rgb]{0.00,0.44,0.13}{#1}}
\newcommand{\PreprocessorTok}[1]{\textcolor[rgb]{0.74,0.48,0.00}{#1}}
\newcommand{\RegionMarkerTok}[1]{#1}
\newcommand{\SpecialCharTok}[1]{\textcolor[rgb]{0.25,0.44,0.63}{#1}}
\newcommand{\SpecialStringTok}[1]{\textcolor[rgb]{0.73,0.40,0.53}{#1}}
\newcommand{\StringTok}[1]{\textcolor[rgb]{0.25,0.44,0.63}{#1}}
\newcommand{\VariableTok}[1]{\textcolor[rgb]{0.10,0.09,0.49}{#1}}
\newcommand{\VerbatimStringTok}[1]{\textcolor[rgb]{0.25,0.44,0.63}{#1}}
\newcommand{\WarningTok}[1]{\textcolor[rgb]{0.38,0.63,0.69}{\textbf{\textit{#1}}}}
\usepackage{longtable,booktabs,array}
\newcounter{none} % for unnumbered tables
\usepackage{calc} % for calculating minipage widths
% Correct order of tables after \paragraph or \subparagraph
\usepackage{etoolbox}
\makeatletter
\patchcmd\longtable{\par}{\if@noskipsec\mbox{}\fi\par}{}{}
\makeatother
% Allow footnotes in longtable head/foot
\IfFileExists{footnotehyper.sty}{\usepackage{footnotehyper}}{\usepackage{footnote}}
\makesavenoteenv{longtable}
\setlength{\emergencystretch}{3em} % prevent overfull lines
\providecommand{\tightlist}{%
  \setlength{\itemsep}{0pt}\setlength{\parskip}{0pt}}
\usepackage{bookmark}
\IfFileExists{xurl.sty}{\usepackage{xurl}}{} % add URL line breaks if available
\urlstyle{same}
\hypersetup{
  hidelinks,
  pdfcreator={LaTeX via pandoc}}

\author{}
\date{}

\begin{document}

\section{LayerForge: Deterministic Layer Decomposition for AI Output
Convergence --- An Empirical Report on Newman Modularity Outperforming
CPM in Small-N Text
Domains}\label{layerforge-deterministic-layer-decomposition-for-ai-output-convergence--an-empirical-report-on-newman-modularity-outperforming-cpm-in-small-n-text-domains}

\begin{quote}
\textbf{Style}: Report framing --- not "we claim this" but "we found
this to hold". Following a publication philosophy of reader-time-cost
\textgreater{} return, we avoid overclaiming and center the paper on
reproducible observations. \textbf{License}: MIT (same as the rest of
the repository). \textbf{Reproducibility}: All raw CSV / scripts / JSON
outputs for the measurements are bundled under
\texttt{scripts/k\_sweep/}, pinned by commit hash.
\end{quote}

\begin{center}\rule{0.5\linewidth}{0.5pt}\end{center}

\subsection{Abstract}\label{abstract}

Conversational AI output tends to include content beyond what the
questioner is actually concerned with. We report on \textbf{LayerForge},
a deterministic pipeline that decomposes a corpus into 4±1 thematic
layers and returns only the layer(s) relevant to a query. The pipeline
combines five established components --- Cowan\textquotesingle s
working-memory upper bound, Newman modularity for community detection,
SCA-style per-layer distillation, HERCULES-style hierarchical KMeans,
and a deterministic core with no LLM in the analysis loop --- chosen so
that "AI noise filtering" does not itself depend on an AI. We
re-implemented a CPM (Constant Potts Model) backend (MIT-pure,
self-implemented) to test whether the theoretically
resolution-limit-free alternative would outperform Newman in this
small-N text domain (operational target N=20-40, external scaling check
up to N=100 on 20 Newsgroups). Across 119-row dual-method correlation
sweep, 196-cell N×K heatmap, 95-row ARI/NMI comparison, and the external
20-Newsgroups benchmark, \textbf{the Newman backend exceeded CPM by a
large effect size in our test range} (H\_struct exact match 87.5\% vs
19\%; 20NG ARI 0.430 vs 0.239 at default K=3, Newman reaches 0.557 when
K=4 is forced; same direction at N=100), against the naïve theoretical
expectation that CPM should win. \textbf{Positioning}: LayerForge shares
its core substrate (sentence-embedding similarity + community detection)
with prior work in clustering-based topic modeling, notably G2T (Zhang
et al. 2023). The contribution of this paper is twofold: (i) the
implementation report of combining SCA decomposition, HERCULES
hierarchical skeleton, the 4±1 cognitive constraint, and a Mode A/B/C
operational layer, and (ii) the systematic empirical comparison of
Newman vs CPM in the small-N text regime, including a three-failure-mode
mechanism analysis (§5). Concurrently with this work, Bohdal et al.
(Samsung 2026) independently arrived at the same problem space on the
LLM memory token side; we treat this convergence as evidence that the
design space is being identified by multiple independent groups. We
report an empirically derived method-selection rule (concept-only;
runtime implementation deferred because in the studied domain the rule
would output "newman" in essentially all cells). The artifact is
published as a Claude Code skill (LayerForge itself does not hold an API
key; LLM access at Boundary 1/2 is mediated through Claude Code), and is
reproducible end-to-end.

\begin{center}\rule{0.5\linewidth}{0.5pt}\end{center}

\subsection{§1. Motivation --- Converging by
concern}\label{1-motivation--converging-by-concern}

Conversation with AI in natural language has a property that looks
"convenient" at first glance but structurally creates a
\textbf{disadvantage}:

\begin{quote}
\textbf{When an AI answers a question, it tends to present related
information comprehensively.} But the user\textquotesingle s "concern"
at that moment in the dialogue is \textbf{usually just one} (e.g., "I
only want to know about the CI/CD of this PR", "I only need to confirm
the assumptions of the experimental design").
\end{quote}

Comprehensive answers create cost in the following three forms:

\begin{enumerate}
\def\labelenumi{\arabic{enumi}.}
\item
  \textbf{Reader cognitive load} The semantic units a human can hold at
  once are bound by short-term memory capacity. Cowan (2010) estimates
  the central working-memory capacity at 3-5 chunks. When AI output
  exceeds this and includes "out-of-concern topics", the reader pays an
  \textbf{additional cognitive cost to extract the parts that match the
  concern}.
\item
  \textbf{Bandwidth waste in multi-agent coordination} When AI agent
  A\textquotesingle s output is consumed downstream by AI agent B,
  comprehensive output from A forces B to either \textbf{pay the cost of
  picking only the concern-relevant region from A\textquotesingle s
  entire output}, or \textbf{drift due to being dragged by
  out-of-concern information}.
\item
  \textbf{Linear waste of AI cost} Processing the full context every
  turn causes token-priced cost to grow linearly. The essential
  information per concern is often smaller than the full context, so
  unnecessary computation is included.
\end{enumerate}

The three are independent phenomena, but \textbf{"AI output not
converging by concern"} is an \textbf{overlapping cause} contributing to
all three (token cost also depends on full-context reprocessing, so the
source is not perfectly identical, but concern-wise convergence is a
partial solution). LayerForge can simultaneously affect all three
through convergence of output.

\subsubsection{1.1 Research question}\label{11-research-question}

\begin{quote}
Can we reduce the three costs above simultaneously by \textbf{converging
the AI\textquotesingle s output into layers that correspond to each
concern}? And can this "layering" be performed \textbf{deterministically
without re-introducing AI}?
\end{quote}

The second question is the important one. \textbf{Filtering AI output
with AI} produces different outputs for the same input (because of the
LLM\textquotesingle s stochastic behavior), so the reproducibility of
concern-wise convergence is lost. This study \textbf{allows AI
involvement only at the input boundary, and implements the filtering
layer itself deterministically}.

Note that "deterministic" in this paper refers to the \textbf{property
of the analysis core (detailed in R3)}, and does not include the LLM
that may optionally intervene at the input boundary (Boundary 1: natural
language → structured) or the output boundary (Boundary 2: structured →
natural language). End-to-end determinism including LLM involvement
holds when the Boundaries are switched to mechanical fallback (see
§3.1).

\subsubsection{1.2 Related work and
positioning}\label{12-related-work-and-positioning}

LayerForge\textquotesingle s design space (sentence embedding +
community detection on similarity graph + per-cluster representation +
clustering-driven context compression) overlaps with several surrounding
research lines, so its position relative to prior art needs to be
stated. Below we organize the main literature in four adjacent groups,
then state LayerForge\textquotesingle s differentiation.

\paragraph{1.2.1 Clustering-based topic modeling and graph-based text
clustering}\label{121-clustering-based-topic-modeling-and-graph-based-text-clustering}

A family that clusters document/passage embeddings to extract topic /
cluster structure:

\begin{itemize}
\tightlist
\item
  \textbf{BERTopic} (Grootendorst 2022, arXiv:2203.05794) ---
  transformer-based document embedding + UMAP dimensionality reduction +
  HDBSCAN clustering + class-based TF-IDF. Per-document
  \textbf{single-topic} assignment is the default.
\item
  \textbf{Top2Vec} (Angelov 2020, arXiv:2008.09470) --- Doc2Vec /
  Sentence Transformer joint embedding + UMAP + HDBSCAN. Topic count is
  determined automatically.
\item
  \textbf{CETopic} (Zhang, Fang, Chen, Namazi-Rad, NAACL 2022, "Is
  Neural Topic Modelling Better than Clustering?") --- SimCSE embedding
  + UMAP + K-means + TF-IDF×IDFi variant, with direct comparison against
  neural topic models.
\item
  \textbf{Vec2GC} (Rao \& Chakraborty 2021, arXiv:2104.09439) ---
  community detection on weighted similarity graphs over terms /
  documents; supports density-based and hierarchical clustering.
\item
  \textbf{G2T (Graph2Topic)} (Zhang, Liu, Yan 2023, arXiv:2304.06653)
  --- pretrained LM embedding → semantic graph construction → community
  detection for topic identification → TFIDF variant for word-topic
  distribution. \textbf{The closest prior art to
  LayerForge\textquotesingle s design substrate.} Claims SOTA against
  BERTopic / CETopic in automatic evaluation.
\item
  \textbf{SCA} (Eichin et al. 2024, arXiv:2410.21054) --- introduces
  \textbf{per-document multi-topic distribution} on top of clustering,
  relaxing the single-assignment assumption of plain clustering. The
  reference implementation for R3 (distillation) in this study.
\item
  \textbf{LLM-Assisted Topic Reduction for BERTopic} (Janssens, Bogaert,
  Van den Poel 2025, arXiv:2509.19365) --- a hybrid family that adds an
  LLM merge step on top of BERTopic output. Differs from the determinism
  principle (R5) of this study.
\end{itemize}

These works adopt the task formulation of "\textbf{classifying
documents/passages into topic groups}". BERTopic / Top2Vec / CETopic
dominantly use density-based (HDBSCAN) or K-means clustering and do not
use community detection; Vec2GC / G2T are in the graph-based
community-detection family. \textbf{G2T is the closest substrate to this
work}, but the community-detection algorithm name (Louvain / CPM /
Leiden) is not stated at the abstract level, and a systematic
Newman-vs-CPM comparison is not performed inside G2T (leaving the
originality of §4-§5 of this work intact).

\paragraph{1.2.2 Hierarchical RAG and
GraphRAG}\label{122-hierarchical-rag-and-graphrag}

A family combining hierarchical structure with retrieval /
summarization:

\begin{itemize}
\tightlist
\item
  \textbf{RAPTOR} (Sarthi et al. 2024, ICLR) --- community detection +
  LLM summarization for hierarchical RAG.
\item
  \textbf{GraphRAG} (Edge et al. 2024, arXiv:2404.16130) --- Leiden +
  LLM summarization for knowledge-graph structuring and query-focused
  summarization.
\item
  \textbf{HippoRAG} --- PageRank-based retrieval, using graph structure
  on the index side.
\item
  \textbf{Core-based Hierarchies for Efficient GraphRAG} (Hossain \&
  Sarıyüce 2026, arXiv:2603.05207) --- replaces
  GraphRAG\textquotesingle s Leiden clustering with \textbf{k-core
  decomposition} to obtain a deterministic hierarchy. The Leiden
  critique is grounded in "modularity optimization admits exponentially
  many near-optimal partitions on sparse graphs", which is \textbf{the
  same problem attacked from a different angle (replacement)} as the
  Good (2010) Q-degeneracy we cite in §5. LayerForge stays inside the
  modularity family and instead documents the failure modes.
\item
  \textbf{RAG vs. GraphRAG: A Systematic Evaluation} (Han et al. 2026,
  arXiv:2502.11371) --- proposes a unified evaluation framework for RAG
  / RAPTOR / HippoRAG2 / GraphRAG. LayerForge sits in the "passage graph
  + decoupled summarization" quadrant of this framework.
\item
  \textbf{GraphRAG-Bench} (Xiao et al. 2025, arXiv:2506.02404) --- a
  domain-specific reasoning benchmark for the GraphRAG family.
\item
  \textbf{HERCULES} (\texttt{bandeerun/pyhercules}, MIT 2025) --- a
  reference implementation of LLM-integrated hierarchical KMeans. No
  peer-reviewed publication exists at the time of this work; we use it
  not as a citation base but as a \textbf{design skeleton reference}
  (§2.4).
\end{itemize}

The bulk of these works assume \textbf{LLM integration in the analysis
pipeline}. LayerForge, by the determinism principle (R5),
\textbf{excludes LLM from the analysis core} and restricts LLM
involvement to Boundary 1/2 (optional) --- an architectural difference.
With Core-based GraphRAG we share the problem framing of "addressing
Leiden\textquotesingle s reproducibility issue", but diverge in the
solution direction: not replacement (k-core), but \textbf{failure-mode
documentation + method-selection signal} inside the modularity family.

\paragraph{1.2.3 Context and memory
compression}\label{123-context-and-memory-compression}

A family compressing input under the LLM\textquotesingle s
context-budget constraint (recent works):

\begin{itemize}
\tightlist
\item
  \textbf{Clustering-driven Memory Compression for On-device LLMs}
  (Bohdal, Saha, Michieli, Ozay, Ceritli 2026, arXiv:2601.17443) ---
  \textbf{the closest concurrent prior art to LayerForge in framing.}
  The mechanism "groups memories by similarity and merges them within
  clusters prior to concatenation" is isomorphic to
  LayerForge\textquotesingle s Mode C (multi-agent bandwidth reduction).
  However, the target is \textbf{user-specific memory tokens} (for
  personalization), and the goal is preserving personalization under a
  context-budget constraint. LayerForge differs on four axes --- (i)
  \textbf{target} (memory tokens vs. corpus passages), (ii)
  \textbf{goal} (personalization vs. concern-wise convergence), (iii)
  \textbf{method} (similarity averaging + merge vs. community detection
  on similarity graph), (iv) \textbf{cognitive constraint} (none vs.
  Cowan 4±1 explicit). Detailed in §1.2.5.
\item
  \textbf{EDU-based Context Compressor} (Zhou et al. 2025,
  arXiv:2512.14244) --- structure-then-select compression via Elementary
  Discourse Unit decomposition. Orthogonal to this work in that the
  compression mechanism is discourse-structural rather than
  clustering-based.
\end{itemize}

This family shares with our work the framing of
"\textbf{clustering-driven context compression}". The fact that our
\textbf{§1 motivation (the three axes of cognitive load / bandwidth /
cost)} and Bohdal et al.\textquotesingle s "balances context efficiency
and personalization quality" arrived at the same problem space
independently means LayerForge is not an isolated design but \textbf{a
position within a problem space that multiple research groups are
working on in parallel in 2025-2026}.

\paragraph{1.2.4 Community detection method
comparisons}\label{124-community-detection-method-comparisons}

A family that systematically benchmarks / theoretically analyzes
multiple community-detection methods:

\begin{itemize}
\tightlist
\item
  \textbf{Lancichinetti \& Fortunato (2009)} --- proposed the LFR
  synthetic benchmark, the standard evaluation infrastructure for
  community-detection algorithms.
\item
  \textbf{Aldecoa \& Marín (2013, Scientific Reports)} ---
  systematically evaluated Newman modularity / CPM / RB / RN / MCL /
  Infomap, etc. on closed benchmarks, reporting that CPM "produces
  accurate results only in the easy ends of the benchmark" and "shows
  some instability". Our §4-§5 small-N text results can be positioned as
  a contrast axis against these synthetic results.
\item
  \textbf{Traag et al. (2011, 2019)} --- the theoretical basis of
  CPM\textquotesingle s resolution-limit-free property (2011) and the
  Leiden refinement (2019).
\item
  \textbf{From Leiden to Pleasure Island: The Constant Potts Model as a
  Hedonic Game} (Felipe, Avrachenkov, Menasche 2025, arXiv:2509.03834)
  --- recasts CPM as a hedonic game (potential game), proves that
  local-utility optimization converges to equilibrium in
  pseudo-polynomial time, and introduces robust partition criteria.
  \textbf{Directly related to §5 of this work (CPM failure-mode
  analysis)}, but on synthetic / random graphs; our small-N text
  similarity graphs are complementary empirical observations.
\end{itemize}

These works center on synthetic graphs / closed benchmarks / theoretical
analysis. To the best of our search, \textbf{a systematic Newman-vs-CPM
comparison on small-N text similarity graphs
(LayerForge\textquotesingle s operational range N=20-40)} does not exist
elsewhere at the time of this work.

\paragraph{1.2.5 LayerForge
positioning}\label{125-layerforge-positioning}

Against the above prior art, this study differentiates on the following
five points:

\begin{enumerate}
\def\labelenumi{\arabic{enumi}.}
\tightlist
\item
  \textbf{vs G2T (Zhang et al. 2023)}: Shares the substrate (sentence
  embedding + community detection on similarity graph), but this study
  focuses on the \textbf{systematic Newman-vs-CPM comparison in the
  small-N text regime + a three-failure-mode classification (§5)}, and
  additionally combines \textbf{SCA distillation (R3), HERCULES skeleton
  (R4), the 4±1 cognitive constraint (R2), and a Mode A/B/C application
  layer}. A direct head-to-head comparison against G2T is out of scope
  for this paper (future work).
\item
  \textbf{vs Clustering-driven Memory Compression (Bohdal et al. 2026,
  Samsung)}: A concurrent prior art that fully matches the top-level
  framing of clustering-driven context compression. We differentiate on
  four axes --- (i) \textbf{target}: LayerForge operates on corpus
  passages, Bohdal on user-specific memory tokens inside the LLM; (ii)
  \textbf{goal}: LayerForge converges AI output by concern (reducing
  cognitive load + bandwidth + cost on three axes), Bohdal fits the
  context budget while preserving personalization; (iii)
  \textbf{method}: LayerForge uses community detection on a similarity
  graph (Newman / CPM switchable), Bohdal uses similarity averaging +
  merge; (iv) \textbf{cognitive constraint}: LayerForge explicitly
  adopts Cowan (2010) 4±1 as an ergonomic upper bound, while Bohdal only
  uses the external context-budget constraint. The two studies are
  \textbf{independent arrivals at the same problem space}, and this
  paper offers a complementary contribution on the corpus-passage side
  by empirically evaluating method choice (Newman vs CPM) inside the
  community-detection family.
\item
  \textbf{vs BERTopic / Top2Vec / CETopic}: This study takes
  \textbf{per-passage hard partitions} as the primary output and
  provides \textbf{deterministic Mode A/B/C} on top of the layer
  structure. BERTopic-family works take per-document topic distributions
  as the primary output, so the task formulation differs and they are
  not suitable as direct baselines (§2.1).
\item
  \textbf{vs RAPTOR / GraphRAG / Core-based GraphRAG / HERCULES}: This
  study \textbf{intentionally decouples LLM from the analysis core} and
  restricts LLM involvement to Boundary 1/2 (R5). This is a design
  choice aimed at \textbf{reproducibility of concern-wise convergence},
  the opposite direction from the LLM-integrated family. With Core-based
  GraphRAG we share the recognition of "Leiden\textquotesingle s
  reproducibility problem", but we diverge by going for
  \textbf{failure-mode documentation + method-selection signal inside
  the modularity family} rather than replacement (k-core).
\item
  \textbf{vs Aldecoa \& Marín / Traag / Felipe et al. (Pleasure
  Island)}: Existing community-detection comparisons center on synthetic
  / random / closed benchmarks. To our knowledge, this paper is the
  first to empirically test "whether resolution-limit-free works" in the
  small-N text domain. The core contribution of §4 / §5 lies here.
\end{enumerate}

In other words, LayerForge is \textbf{an application that places SCA +
HERCULES + 4±1 + Mode A/B/C on top of G2T\textquotesingle s design
space}, and the core contribution narrows to \textbf{(i) the
implementation report of that combination, and (ii) the systematic
Newman-vs-CPM comparison in the small-N text regime + the
three-failure-mode classification (§5)}. We do not claim
"community-detection method comparison itself" as a novel contribution;
the paper should be read as "\textbf{empirical evaluation and
documentation of known methods within this domain}". The framing overlap
with Bohdal et al. (Samsung 2026) does not negate this
work\textquotesingle s novelty; we record it as the fact that
\textbf{two studies independently identified the same problem space}.

\begin{center}\rule{0.5\linewidth}{0.5pt}\end{center}

\subsection{§2. Why these technologies --- Necessity of the
five-component
combination}\label{2-why-these-technologies--necessity-of-the-five-component-combination}

Extracting requirements from the §1 motivation:

{\def\LTcaptype{none} % do not increment counter
\begin{longtable}[]{@{}lll@{}}
\toprule\noalign{}
Requirement & Technology choice & Related literature \\
\midrule\noalign{}
\endhead
\bottomrule\noalign{}
\endlastfoot
(R1) \textbf{Extract "concern-wise units" from a corpus} & Community
detection on similarity graph & Newman (2006), Fortunato \& Barthélemy
(2007), Traag et al. (2011) \\
(R2) \textbf{An ergonomic upper bound for the reader\textquotesingle s
cognition} & 4±1 layers (Cowan\textquotesingle s working memory
capacity) & Cowan (2010), Miller (1956) \\
(R3) \textbf{Extract the "core information" of each layer} & Per-layer
distillation (SCA-style) & Eichin et al. (2024) --- Semantic Component
Analysis \\
(R4) \textbf{Scalability to large corpora} & Hierarchical KMeans
foundation (HERCULES-style) & Reference impl
\texttt{bandeerun/pyhercules} (MIT, 2025) \\
(R5) \textbf{Determinism that "does not filter AI output with AI"} &
LLM-free analysis core, AI involvement only at boundaries & LayerForge
design rationale \\
\end{longtable}
}

\subsubsection{2.1 (R1) Why community
detection}\label{21-r1-why-community-detection}

An established method for \textbf{unsupervised} automatic extraction of
"related sentence groups". RAG extensions (RAPTOR, GraphRAG, HippoRAG)
use Louvain/Leiden for passage clustering under the same spirit. This
work belongs to that family, but differs by \textbf{decoupling community
detection from LLM summary} under the R5 constraint (discussed later).

Methods that output document-level topic distributions (the
topic-modeling family, e.g. LDA, BERTopic) are \textbf{not primarily
aimed at producing passage-level hard partitions}, so they are not
appropriate as direct baselines for this work\textquotesingle s task
formulation of "which layer does each passage belong to". The comparison
target of this paper is restricted to \textbf{choice inside the
community-detection method (Newman vs CPM)}.

\subsubsection{2.2 (R2) Why 4±1}\label{22-r2-why-41}

Cowan (2010) estimates the central short-term-memory capacity at 3-5
chunks. If we require that "concern-wise convergence be easy for the
reader", placing the \textbf{ergonomic upper bound} of layer count here
is natural. This is a design decision to \textbf{adopt this as a target
value}, not a claim that "4±1 is a universal cognitive constant".

\subsubsection{2.3 (R3) Why SCA
distillation}\label{23-r3-why-sca-distillation}

Community detection produces "which sentence belongs to which layer",
but does not produce a \textbf{representative expression} for the layer.
To hand a layer to a reader or downstream agent, a short form of "what
topic this layer is about" is needed. Eichin et al.
(2024)\textquotesingle s Semantic Component Analysis distills
per-cluster representative components, matching this requirement.

\subsubsection{2.4 (R4) Why HERCULES}\label{24-r4-why-hercules}

Flat 1-level community detection is sufficient for the main results of
this work, but we adopt a \textbf{hierarchical KMeans skeleton} in
advance, anticipating future operation such as "recursively expanding an
N=10K large corpus into 4×4×4×4 = 256 levels" (the current
implementation defaults to max\_depth=1 and is recursively expandable
via an option).

Note that HERCULES, at the time of this work, exists \textbf{only as a
GitHub reference implementation (\texttt{bandeerun/pyhercules}, MIT)
with no peer-reviewed publication}. We reference HERCULES as a "prior
OSS implementation that adopted hierarchical KMeans" for positioning,
not as a citation base for its own theoretical superiority. The
substantive basis for R4 is the design judgment that "recursive KMeans
functions as a skeleton for large-corpus extension"; HERCULES is one
implementation example (= reference target) of that.

\subsubsection{2.5 (R5) Why a deterministic
core}\label{25-r5-why-a-deterministic-core}

As stated in §1, \textbf{the stochastic behavior of LLMs destroys the
reproducibility of concern-wise convergence}. If different layerings
emerge for the same input, the "converged output" cannot be trusted.
This work uses:

\begin{itemize}
\tightlist
\item
  \textbf{Input boundary}: natural language → structured (Boundary 1),
  LLM involvement allowed
\item
  \textbf{Analysis core}: deterministic (community detection + KMeans +
  SCA are all deterministic algorithms)
\item
  \textbf{Output boundary}: structured → natural language (Boundary 2),
  LLM involvement allowed (optional; in this work Modes A and C operate
  without AI)
\end{itemize}

Under this constraint, the analysis results guarantee \textbf{same input
+ same hyperparameters → same output}. Here, hyperparameters include the
community-detection method choice
(\texttt{community\_method="newman"\ \textbar{}\ "cpm"}), random seed,
\texttt{target\_range}, embedding model, etc. Reproducibility holds with
these fixed; switching methods can yield different partitions
(demonstrated in §4.3).

\subsubsection{2.6 Summary --- Necessity of the
combination}\label{26-summary--necessity-of-the-combination}

The five technology choices are each derived from independent
requirements (R1-R5), and \textbf{conversely, alternative choices that
satisfy each requirement do not fit this study\textquotesingle s task}:

\begin{itemize}
\tightlist
\item
  Skip community detection → cannot automatically extract concern-wise
  units
\item
  Drop 4±1 → lose the basis for cognitive-load reduction
\item
  Skip SCA → no representative expression per layer
\item
  Drop HERCULES → no base for large-scale expansion
\item
  Drop determinism → self-contradiction by trying to filter AI with AI
  for the §1 problem
\end{itemize}

§3 shows the \textbf{design rationale} for how this combination
implements the §1 motivation (R1-R5). §4 empirically tests the validity
of the \textbf{most important choice inside the combination (the
community-detection method, Newman vs CPM)}. \textbf{Individual
ablations of SCA distillation / HERCULES hierarchy / the 4±1 constraint
/ determinism are out of scope} for this paper, and are stated in §7.5.
The claim of this paper is restricted to "an empirical method comparison
on the most important choice inside the combination, showing that the
Newman backend fits this domain".

\begin{center}\rule{0.5\linewidth}{0.5pt}\end{center}

\subsection{§3. Approach --- LayerForge
implementation}\label{3-approach--layerforge-implementation}

\subsubsection{3.1 Pipeline}\label{31-pipeline}

\begin{verbatim}
[input: list of passages, optional query]
       │
       ▼
(B1) parse_to_structure       — natural language → FormulationInput (optional LLM, has mechanical fallback)
       │
       ▼
build_similarity_matrix       — sentence-transformers cosine similarity
       │
       ▼
find_valid_scale / find_cpm_resolution
                              — find K in the 4±1 range via threshold θ (Newman) or γ (CPM)
       │
       ▼
detect_communities            — Newman spectral + KMeans / CPM-Louvain (community_method option)
       │
       ▼
compute_modularity            — Q (both methods) + cpm_h (when CPM)
       │
       ▼
distill_layer                 — per-layer SCA (UMAP + HDBSCAN-style components)
       │
       ▼
hierarchy_to_layer_summaries  — { layers: [...], inter_layer_relations: [...] }
       │
       ▼
(B2) render_to_natural        — structured → natural language (optional LLM)
       │
       ▼
[output: 4±1 layers + per-layer reps + optional natural-language rendering]
\end{verbatim}

\subsubsection{3.2 Three modes of
operation}\label{32-three-modes-of-operation}

{\def\LTcaptype{none} % do not increment counter
\begin{longtable}[]{@{}lll@{}}
\toprule\noalign{}
Mode & CLI & Purpose \\
\midrule\noalign{}
\endhead
\bottomrule\noalign{}
\endlastfoot
\textbf{A (decompose)} & \texttt{layerforge-decompose} & Decompose a
corpus into 4±1 layers; pull only the layer(s) of interest from the
UI \\
\textbf{B (decide)} & \texttt{layerforge-decide} & Layer-organize
decision information; track concern transitions via
\texttt{open}/\texttt{close}/\texttt{settle} \\
\textbf{C (compress)} & \texttt{layerforge-compress} & Compress an
AI\textquotesingle s verbose output into the layer subset relevant to
the query (decision-less subset guarantee) \\
\end{longtable}
}

\subsubsection{3.3 Claude Code skill
form}\label{33-claude-code-skill-form}

LayerForge runs standalone as a Python CLI, but is operationally placed
as a \textbf{Claude Code skill}:

\begin{itemize}
\tightlist
\item
  \textbf{LayerForge itself holds no API key} (when LLM involvement is
  needed at Boundary 1/2, it is executed via Claude Code, where API keys
  are managed)
\item
  A PostToolUse hook automatically validates the output schema
\item
  \texttt{.claude/skills/layerforge/SKILL.md} functions as \textbf{an
  instruction sheet for the Claude using the skill}
\end{itemize}

In this form, LayerForge can be inserted into existing AI workflows as a
"\textbf{peripheral tool of the LLM}", while \textbf{the analysis core
preserves LLM-free determinism} (Boundary 1/2 allow optional LLM
involvement as per §3.1, with mechanical fallback available).

\subsubsection{3.4 Dual community-detection
backend}\label{34-dual-community-detection-backend}

The central engineering choice of this work is the community-detection
method. The implementation provides both:

\begin{itemize}
\tightlist
\item
  \textbf{Newman (default)}: threshold θ + spectral algorithm + KMeans
  on embeddings; quality measured by modularity Q
\item
  \textbf{CPM (opt-in)}: Leiden-CPM (self-implemented, MIT-pure);
  maximizes the H function with the \texttt{(n\_c\ choose\ 2)} quadratic
  penalty
\end{itemize}

Switching is done via the \texttt{community\_method} option:

\begin{Shaded}
\begin{Highlighting}[]
\NormalTok{layerforge\_core(}\BuiltInTok{input}\NormalTok{, community\_method}\OperatorTok{=}\StringTok{"newman"}\NormalTok{)  }\CommentTok{\# default}
\NormalTok{layerforge\_core(}\BuiltInTok{input}\NormalTok{, community\_method}\OperatorTok{=}\StringTok{"cpm"}\NormalTok{)     }\CommentTok{\# opt{-}in}
\end{Highlighting}
\end{Shaded}

\subsubsection{3.5 Reproducibility
infrastructure}\label{35-reproducibility-infrastructure}

{\def\LTcaptype{none} % do not increment counter
\begin{longtable}[]{@{}ll@{}}
\toprule\noalign{}
Artifact & Purpose \\
\midrule\noalign{}
\endhead
\bottomrule\noalign{}
\endlastfoot
\texttt{scripts/k\_sweep/correlation\_data.py} (119 rows) & sweep over 5
configs × 12 K × 2 methods \\
\texttt{scripts/k\_sweep/heatmap\_N\_x\_K.py} (196 attempted cells = 7 N
× 14 K × 2 methods, of which 167 succeeded --- the remaining 29 were
skipped by the K ≥ N filter or by CPM γ-bisection failure) & matrix over
7 N × 14 K × 2 methods \\
\texttt{scripts/k\_sweep/cpm\_compare.py} (95 rows + ARI/NMI) & direct
Newman-vs-CPM comparison \\
\texttt{scripts/k\_sweep/multi\_corpus\_verify\_v2.py} (8 conditions) &
test whether K\_optimal tracks n\_themes \\
\texttt{scripts/k\_sweep/k10\_multi\_corpus.py} & test whether K=10
self-routing is preserved \\
\texttt{scripts/k\_sweep/run\_robustness.py} (32 setting-method
aggregates, each holding 8 K-range measurements; 256 K-range cells
total) & robustness over 16 settings × 8 K ranges × 2 methods \\
\texttt{tests/integration/test\_real\_data\_20ng.py} & 20 Newsgroups
external benchmark (both methods) \\
\texttt{tests/axioms/test\_cpm\_karate\_club.py} &
Zachary\textquotesingle s Karate Club reference (CPM correctness
gate) \\
\end{longtable}
}

All raw CSV / JSON / PNG are bundled in git, pinned by commit hash.

\begin{center}\rule{0.5\linewidth}{0.5pt}\end{center}

\subsection{§4. Empirical findings --- Observations that
held}\label{4-empirical-findings--observations-that-held}

\begin{quote}
This section is written as a report of "what was found to hold". We
avoid claims such as "superior" or "best", and place \textbf{the
measured numbers and their ranges} at the center.
\end{quote}

\subsubsection{4.0 Terminology for K (preliminaries for this
section)}\label{40-terminology-for-k-preliminaries-for-this-section}

{\def\LTcaptype{none} % do not increment counter
\begin{longtable}[]{@{}lll@{}}
\toprule\noalign{}
Term & Definition & Used in \\
\midrule\noalign{}
\endhead
\bottomrule\noalign{}
\endlastfoot
\texttt{target\_range} & the range of K to search (e.g., 4±1 = (3, 5)) &
argument of \texttt{find\_valid\_scale} / \texttt{find\_cpm\_resolution}
in §3.1 \\
\texttt{K\_actual} & the number of communities obtained as a run result
& all figures / tables \\
\texttt{K\_optimal} & the K that maximizes Q in a K sweep (Newman) &
§4.1, §4.4 \\
\texttt{Q\ peak\ K} & synonym for \texttt{K\_optimal} & §4.1 \\
\texttt{n\_themes} & the ground-truth theme count of the corpus & §4.4
(H\_struct), §4.6 (20NG) \\
\texttt{K=10} & a specific operational candidate value (AI input
compression) & §4.5 \\
\texttt{default\ K} & the probe value in the method-selection rule (K=4
in this paper) & §5.7 \\
\end{longtable}
}

§4.1 sweeps Q peak without pinning \texttt{target\_range} to
\texttt{4±1}. §4.4 measures, for each setting, whether
\texttt{K\_actual} tracks the corresponding \texttt{n\_themes}. The
difference in corpus setups between the two is stated explicitly in
§4.4.

\subsubsection{4.1 N-dependent behavior of Q peak K
(Newman)}\label{41-n-dependent-behavior-of-q-peak-k-newman}

The corpus in this sub-section is a \textbf{cross-domain mpnet corpus}
(real-world markdown corpus of the KDF-perovskite project, four themes
--- philosophy / exploration / proof / blog, \texttt{n\_themes\ =\ 4}),
sampled at \texttt{per\_theme\ =\ 2,3,4,5,6,8,10} to vary N (a different
corpus family from the synthetic disjoint-vocabulary corpus of §4.4; see
§4.4\textquotesingle s reconciliation for the relation). The measurement
question is \textbf{"does Q peak K move when the sampling density of the
same corpus structure varies?"} --- a check of whether the theoretical Q
degeneracy of Good (2010) is reproduced in this implementation.

Q max over 7 N values × 14 K values × Newman backend:

{\def\LTcaptype{none} % do not increment counter
\begin{longtable}[]{@{}rrr@{}}
\toprule\noalign{}
N & Q peak K & Q peak value \\
\midrule\noalign{}
\endhead
\bottomrule\noalign{}
\endlastfoot
8 & 3 & 0.403 \\
12 & 6 & 0.778 \\
16 & 11 & 0.720 \\
20 & 6 & 0.716 \\
24 & 9 & 0.646 \\
32 & 4 & 0.550 \\
40 & 6 & 0.610 \\
\end{longtable}
}

\textbf{Observation}: For the same corpus structure, varying N (sample
size) makes Q peak K fluctuate in the range 3-11. This demonstrates that
the Newman Q degeneracy theoretically established by Good et al. (2010)
is reproduced in this implementation. The independent contribution of
this paper is not the discovery itself, but \textbf{the empirical
dataset on the N×K axes}.

\subsubsection{4.2 Method-agnostic monotonicity of the above-limit
fraction}\label{42-method-agnostic-monotonicity-of-the-above-limit-fraction}

Using the Fortunato \& Barthélemy (2007) resolution limit (√(L/2)) as a
baseline, we compute the fraction of communities exceeding the limit. We
present as the primary result the \textbf{common-reference measurement
that directly supports the method-agnostic claim}
(\texttt{scripts/k\_sweep/data\_current/cpm\_compare\_data.csv},
comparing both methods on the same graph thresholded by the corpus-wide
median similarity):

{\def\LTcaptype{none} % do not increment counter
\begin{longtable}[]{@{}rrrr@{}}
\toprule\noalign{}
N & K & Newman above-limit & CPM above-limit \\
\midrule\noalign{}
\endhead
\bottomrule\noalign{}
\endlastfoot
12 & 3 & 0.67 & 0.67 \\
12 & 4 & 0.25 & 0.25 \\
20 & 4 & 0.50 & 0.50 \\
24 & 4 & 0.75 & 0.75 \\
24 & 5 & 0.60 & 0.60 \\
40 & 5 & 0.80 & 0.60 \\
\end{longtable}
}

\textbf{Observation}: In 5 of 6 cells the above-limit fraction matches
\textbf{exactly} between the two methods; only N=40, K=5 diverges
(Newman 0.80 / CPM 0.60). Both Newman and CPM are monotone-decreasing in
K, and the two curves are nearly the same shape in this
paper\textquotesingle s operational range (N=12-40). The N=40 divergence
is consistent with the regime §4.3 shows ("partitions become divergent
at N≥32") --- on the same graph, different methods produce different
partition shapes and hence different cluster-size distributions, which
can move the above-limit fraction apart.

→ This work empirically shows (within the tested range) that \textbf{the
above-limit fraction functions as a supplementary K-selection signal
that does not depend on the community-detection method}. The
asymmetric-reference version of the same cells, measured at the θ chosen
by Newman\textquotesingle s own \texttt{find\_valid\_scale}, is repeated
in §5.2 and handled on the mechanism-analysis side (to separate the
insights from the two reference choices).

\subsubsection{4.3 Newman vs CPM partition agreement (ARI by
N)}\label{43-newman-vs-cpm-partition-agreement-ari-by-n}

Running both methods on the same (N, K) cell and measuring the Adjusted
Rand Index between the partitions:

{\def\LTcaptype{none} % do not increment counter
\begin{longtable}[]{@{}rrr@{}}
\toprule\noalign{}
N & ARI mean & ARI max \\
\midrule\noalign{}
\endhead
\bottomrule\noalign{}
\endlastfoot
12 & 0.871 & 1.000 \\
20 & 0.728 & 1.000 \\
24 & 0.689 & 1.000 \\
32 & 0.490 & 0.779 \\
40 & 0.411 & 0.842 \\
\end{longtable}
}

\textbf{Observation}: At N≤24 the partitions nearly coincide (perfect
match at some K), while at N≥32 they diverge (at most 0.85). The naïve
expectation of "Newman and CPM produce the same result" is \textbf{valid
only on the small-N side} of this domain; in the operational range
(N=20-40), method choice was found to affect partition structure.

\subsubsection{4.4 H\_struct (Q peak K = n\_themes)
tracking}\label{44-h_struct-q-peak-k--n_themes-tracking}

This sub-section measures a \textbf{different axis} from §4.1:

\begin{itemize}
\tightlist
\item
  §4.1: "how does Q peak K move \textbf{when N varies for the same
  corpus structure}" (Q degeneracy reproduction)
\item
  §4.4: "for \textbf{corpora with different n\_themes}, does Q peak K
  match n\_themes" (test of the H\_struct hypothesis)
\end{itemize}

16 settings = \textbf{4 n\_themes (3, 4, 5, 7) × 2 embedders (MiniLM,
mpnet) × 2 seeds (42, 123)}. For each setting we sweep K ranges over
\texttt{\{1-2,\ 2-3,\ 3-4,\ 3-5,\ 5-7,\ 6-8,\ 8-10,\ 10-12\}}. The
two-tier design intent is: (a) continuously covering K=1 through 12 with
overlap to confirm binary-search stability of
\texttt{find\_valid\_scale} over the whole K span; (b) densely covering
Cowan\textquotesingle s 4±1 (=3-5) with three ranges
(\texttt{2-3,\ 3-4,\ 3-5}) to reduce error in this
paper\textquotesingle s region of interest (4±1). For each range, we
take the K\_actual that maximizes Q and compare against the
corresponding \texttt{n\_themes}:

\begin{itemize}
\tightlist
\item
  Newman: \textbf{14/16 (87.5\%) exact match}
  (\texttt{K\_actual\ ==\ n\_themes})
\item
  CPM: \textbf{3/16 (19\%) exact match}
\end{itemize}

\textbf{Observation}: The H\_struct hypothesis "Q peak K tracks the
corpus\textquotesingle s natural theme count" is \textbf{strongly
supported for the Newman backend in this domain of 16 settings
(synthetic disjoint-vocabulary corpus)}, and not supported for the CPM
backend. Since 16 settings is a small sample, we do not perform formal
statistical significance tests (e.g., binomial test) in this paper, and
qualitatively evaluate the effect size (87.5\% vs 19\%).

\textbf{Scope limitation (important)}: The 87.5\% tracking above is
\textbf{a result on a synthetic disjoint-vocabulary corpus}, and this
paper does \textbf{not} make the strong claim that
"Newman\textquotesingle s \texttt{find\_valid\_scale} pipeline
automatically tracks n\_themes on any corpus". On the external benchmark
(20 Newsgroups, N=100, §4.6), even Newman returns K=3 under the default
\texttt{target\_range=(3,5)} and fails to track n\_themes=4. This
reflects the gap between the synthetic disjoint-vocabulary condition and
public-corpus vocabulary overlap, and is positioned mechanistically in
§5.4 (the \texttt{find\_cpm\_resolution}-style calibration design
prefers the lower end of \texttt{target\_range}). The H\_struct claim of
this section should be read as \textbf{partial support under the
synthetic-corpus condition}.

→ \textbf{Figure}: \texttt{figures/new/fig\_h\_struct.png} --- a 2-panel
scatter (Newman \textbar{} CPM) of K\_actual vs n\_themes, with a
diagonal indicating perfect tracking. The Newman side concentrates near
the diagonal, while the CPM side systematically falls below (visualizing
under-merging).

\textbf{Relation to §4.1 (reconciliation)}: §4.1 and §4.4 use different
corpus families:

{\def\LTcaptype{none} % do not increment counter
\begin{longtable}[]{@{}lll@{}}
\toprule\noalign{}
Item & §4.1 & §4.4 \\
\midrule\noalign{}
\endhead
\bottomrule\noalign{}
\endlastfoot
Corpus origin & \textbf{real-world cross-domain} (philosophy /
exploration / proof / blog markdown of the KDF-perovskite project) &
\textbf{synthetic disjoint-vocabulary}
(\texttt{scripts/k\_sweep/corpora.py::make\_corpus}, each theme using
unique entity names + property templates) \\
Axis varied & sampling density (per\_theme) for the same corpus
structure & different n\_themes (3,4,5,7) × embedders × seeds \\
Primary measurement & fluctuation of Q peak K with N & match between Q
peak K and n\_themes for each setting \\
Result & Q peak K fluctuates in 3-11 with N (Q degeneracy reproduction)
& 14/16 (87.5\%) Q peak K = n\_themes \\
\end{longtable}
}

The two measure \textbf{different questions on different corpus
families} and do not contradict: §4.1 measures the "degenerate effect of
sampling density inside the same corpus" on a real-world corpus, and
§4.4 measures "tracking across corpora with different n\_themes" on
synthetic corpora. The real-world cross-domain corpus has more
vocabulary overlap than §4.4\textquotesingle s synthetic corpus, so it
is expected behavior that Q peak K in §4.1 does not pin to n\_themes=4.

\subsubsection{4.5 K=10 self-routing (AI input compression
candidate)}\label{45-k10-self-routing-ai-input-compression-candidate}

\textbf{Definition of self-routing accuracy}: for each passage
\texttt{p}, treat the text of \texttt{p} itself as a query, compare the
embedding of \texttt{p} against all layer centroids, and select the
layer with the closest centroid. The fraction of cases in which the
chosen layer matches \texttt{p}\textquotesingle s layer assignment. This
is a \textbf{weaker test} than LLM-side paraphrased query routing, but
sufficient to verify algorithm-level partition consistency. Production
query routing is out of scope for this paper (see §7 Limitations).

8 conditions (2 corpora × 2 embedders × 2 methods) × K=10 cell:

{\def\LTcaptype{none} % do not increment counter
\begin{longtable}[]{@{}lrr@{}}
\toprule\noalign{}
condition & Newman self-routing & CPM self-routing \\
\midrule\noalign{}
\endhead
\bottomrule\noalign{}
\endlastfoot
same-domain 5themes / MiniLM & 30/30 (100\%) & 29/30 (97\%) \\
cross-domain 4themes / MiniLM & 24/24 (100\%) & 23/24 (96\%) \\
same-domain 5themes / mpnet & 30/30 (100\%) & 30/30 (100\%) \\
cross-domain 4themes / mpnet & 24/24 (100\%) & 24/24 (100\%) \\
\end{longtable}
}

\textbf{Observation}: At K=10 self-routing stays in 96-100\%, and the
theoretical minimum read-fraction on the partition-consistency metric is
1/K = 10\% (a 10× reduction). The two methods have similar accuracy.
\textbf{On the self-routing metric, K=10 as an AI-input compression
target holds method-agnostically.}

\textbf{Important scope limitation}: The "10× reduction" above is the
\textbf{theoretical minimum value of partition consistency}, and
\textbf{production retrieval accuracy under paraphrased queries is not
measured in this section}. What self-routing measures is intra-cluster
tightness --- "does a passage, used as its own query, return to the
layer it belongs to?" --- and there is a gap between this and production
query routing (which involves paraphrase and query-passage semantic
gap). Reading these numbers as "achieving 10× AI input compression"
would be a leap; this section reports only that the \textbf{necessary
condition for partition consistency} is satisfied. Verification of
production query routing accuracy is out of scope for this paper (see
§7.2 / §7.5).

\textbf{Granularity dependence vs §4.4 / §4.6}: The phenomenon that the
method-difference size is very different between this section (K=10) and
§4.4 (K = n\_themes) is explained mechanistically in §5. In short:
\textbf{in the K = n\_themes regime, Newman picks up the true thematic
structure via the Q peak while CPM under-merges}, so method differences
become visible; meanwhile \textbf{in over-clustering regimes such as
K=10, both methods are forced into "more clusters than themes"}, so
partition fragmentation dominates over the ground-truth distinction and
method differences become harder to observe. The granularity regime
determines whether method differences surface.

\subsubsection{4.6 20 Newsgroups (external
benchmark)}\label{46-20-newsgroups-external-benchmark}

Four thematically distinct newsgroups (sci.med / sci.space /
rec.sport.hockey / talk.politics.guns) with 25 docs/topic = \textbf{N =
100 documents} (only in this sub-section, serving also as a scaling
check from N=20-40 in §4.1-4.5).

\textbf{Measured K}: both methods are run with
\texttt{target\_range\ =\ (3,\ 5)} (Cowan\textquotesingle s 4±1 default
range). In the current default behavior of the implementation, both
methods return K=3 (\texttt{tests/integration/test\_real\_data\_20ng.py}
admits Newman \{3,4,5\} and CPM \{2,3,4,5\}). \textbf{The phenomenon
that both methods plateau at K=3 is a symptom of the K-calibration
design of this implementation (both \texttt{find\_valid\_scale} and
\texttt{find\_cpm\_resolution} are specified to prefer the lower end of
\texttt{target\_range}; mechanism explained in §5.4), and the failure to
reach n\_themes=4 is not an intrinsic limit of Newman.} The Newman ARI
of 0.557 under forced K=4 shows the level reachable without going
through calibration. The obtained partition is compared against the
\textbf{ground truth (n\_themes = 4)} and Adjusted Rand Index is
measured (sklearn.metrics.adjusted\_rand\_score):

{\def\LTcaptype{none} % do not increment counter
\begin{longtable}[]{@{}lrr@{}}
\toprule\noalign{}
method & K\_actual (default) & ARI vs ground truth \\
\midrule\noalign{}
\endhead
\bottomrule\noalign{}
\endlastfoot
Newman & 3 & \textbf{0.430} \\
CPM & 3 & \textbf{0.239} \\
chance baseline (random partition) & --- & ≈ 0 \\
\end{longtable}
}

For reference, fixing K=4 yields Newman ARI \textbf{0.557}, and K=5
yields 0.313. CPM stays at ARI ≈ 0.24 across K=3,4,5 (under-merging
keeps it in the same partition family regardless of K).

→ \textbf{Figure}: \texttt{figures/new/fig\_20ng\_ari.png} --- bar chart
(Newman / CPM × K=3,4,5), with the chance-baseline line and
Newman\textquotesingle s K=4 best (0.557) emphasized in annotations. The
gap of Newman 1.8× at default K=3 and Newman 2.3× under forced K=4 is
visually clear.

\textbf{Observation}: The external validity on a public benchmark is
that \textbf{at the default setting, the Newman backend exceeds the CPM
backend by about a factor of 1.8} (Newman 0.430 / CPM 0.239). Both
methods clearly exceed the chance baseline, and Newman\textquotesingle s
advantage is observed stably. Newman\textquotesingle s ARI=0.557 under
forced K=4 shows that Newman can reach higher ARI at an appropriate K.
CPM plateaus at 0.24 regardless of K, consistent with the (n\_c choose
2) penalty effect (mode (c)) of §5.3.

\begin{center}\rule{0.5\linewidth}{0.5pt}\end{center}

\subsection{§5. Why CPM Underperforms --- Three failure modes and
operational
policy}\label{5-why-cpm-underperforms--three-failure-modes-and-operational-policy}

The most counter-intuitive result observed in §4 is that \textbf{CPM is
consistently outperformed by Newman}. Traag et al. (2011) mathematically
established CPM\textquotesingle s resolution-limit-freeness via the
subgraph-invariance property, positioning it as a \textbf{correct fix}
for Newman\textquotesingle s Q degeneracy (Good 2010). Nevertheless, in
this study\textquotesingle s small-N text domain CPM degraded in the
opposite direction. This section organizes the mechanism by separating
it into \textbf{three distinct failure modes}. The opposite-direction
behaviors of K=4 over-split on Karate Club (§7.3) and K=3 under-merge on
20 Newsgroups (§4.6) are positioned consistently as different modes.

\subsubsection{5.1 The three failure
modes}\label{51-the-three-failure-modes}

The causes of CPM-Louvain\textquotesingle s deviation from the ground
truth can be separated, within the verification range of this work, into
\textbf{at least three independent mechanisms}:

{\def\LTcaptype{none} % do not increment counter
\begin{longtable}[]{@{}llll@{}}
\toprule\noalign{}
Mode & Content & Where it mainly appears & Section in this paper \\
\midrule\noalign{}
\endhead
\bottomrule\noalign{}
\endlastfoot
(a) \textbf{Louvain refinement gap} & There exist macro partitions
unreachable by vanilla Louvain\textquotesingle s single-node moves (a
known problem that Traag 2019 fixed with Leiden). Cannot escape even by
lowering γ. & dense small graphs (e.g., Karate Club), ground truth being
a macro split & §5.5 / §7.3 \\
(b) \textbf{Calibration bias} & This implementation\textquotesingle s
\texttt{find\_cpm\_resolution} bisects γ and is designed to
\textbf{prefer the smallest K within target\_range}
(\texttt{cpm\_backend.py::find\_cpm\_resolution} L262-263). If n\_themes
is at the upper end of target\_range, the result converges to the lower
K and under-merges. & external corpora where n\_themes does not match
target\_range{[}0{]} & §5.4 \\
(c) \textbf{(n\_c choose 2) penalty effect} & In regimes where the
resolution limit does not bind under small N, CPM\textquotesingle s
quadratic penalty relatively dominates \texttt{m\_c}, selecting coarser
partitions. ARI degradation against Newman remains even at fixed K. &
small-N text domain (this paper\textquotesingle s operational range),
K-matched comparison & §5.3 \\
\end{longtable}
}

Attribution of the numerical gaps in this paper:

\begin{itemize}
\tightlist
\item
  \textbf{§4.4 H\_struct 87.5\% vs 19\%} (synthetic disjoint-vocabulary,
  K\_actual taken at Q peak for each setting) → primarily (c). The K
  range densely covers around n\_themes, so the influence of (b) is
  small.
\item
  \textbf{§4.6 default K=3 (20NG)} → dominated by (b) calibration bias.
  Both methods converge to the lower end K=3 of target\_range=(3,5).
\item
  \textbf{§4.6 Newman 0.557 vs CPM 0.24 under forced K=4} → (c) penalty
  effect. The residual CPM underperformance even after avoiding
  calibration.
\item
  \textbf{§7.3 Karate K=4 over-split} → (a) Louvain refinement gap.
  Cannot reach the K=2 macro split regardless of γ.
\end{itemize}

\subsubsection{5.2 Structural difference of the objective
functions}\label{52-structural-difference-of-the-objective-functions}

Comparing the objective functions of the two methods:

\begin{itemize}
\tightlist
\item
  Newman:
  \texttt{Q\ =\ (1/2L)\ Σ\_ij\ {[}A\_ij\ -\ (k\_i\ k\_j\ /\ 2L){]}\ δ(c\_i,\ c\_j)}

  \begin{itemize}
  \tightlist
  \item
    Internal edges compared against a null model (expectation k\_i k\_j
    / 2L)
  \item
    No explicit penalty on cluster size
  \end{itemize}
\item
  CPM: \texttt{H\ =\ Σ\_c\ {[}m\_c\ -\ γ\ ×\ (n\_c\ choose\ 2){]}}

  \begin{itemize}
  \tightlist
  \item
    From intra-edges, \textbf{the cluster pair count × γ} is subtracted
  \item
    Penalty grows \textbf{quadratically} in the cluster size n\_c
  \end{itemize}
\end{itemize}

→ \textbf{Figure}: \texttt{figures/new/fig\_cpm\_mechanism.png} --- a
comparison plot of penalty/null-model contribution as a function of
cluster size n\_c. The CPM \texttt{(n\_c\ choose\ 2)} quadratic curve
exceeds Newman\textquotesingle s linear null-model reference, with the
"under-merging zone" shaded.

The \textbf{common-reference} measurement shown in §4.2 (both methods
compared on the same graph thresholded at the corpus\textquotesingle s
median similarity) has above-limit fraction matching exactly in 5 of 6
cells in the operational range N=12-40. That is, in this domain
\textbf{Newman\textquotesingle s resolution limit is not operationally
binding}, and the theoretical motivation for adopting CPM
(limit-freeness) targets \textbf{a problem that is not a problem in this
domain}. What remains is only the effect of CPM\textquotesingle s own
penalty structure --- this is the starting observation for §5.3 onward.

For reference, we repeat the same cells under the \textbf{asymmetric
reference} (Newman thresholded at the θ chosen by its own
\texttt{find\_valid\_scale}, CPM thresholded at the median similarity),
which complements §4.2. This is useful as a reference for directly
inspecting "whether Newman\textquotesingle s resolution limit is
binding", but we use the common-reference of §4.2 as the ground for
method comparison:

{\def\LTcaptype{none} % do not increment counter
\begin{longtable}[]{@{}rrrr@{}}
\toprule\noalign{}
N & K & Newman above-limit (asym) & CPM above-limit (asym) \\
\midrule\noalign{}
\endhead
\bottomrule\noalign{}
\endlastfoot
12 & 3 & 0.67 & 0.67 \\
12 & 4 & 0.25 & 0.25 \\
20 & 4 & 0.75 & 0.50 \\
24 & 4 & 0.75 & 0.75 \\
24 & 5 & 0.40 & 0.60 \\
40 & 5 & 0.60 & 0.60 \\
\end{longtable}
}

(The difference on the Newman side comes from the θ-selection path, not
from the partition structure on the graph. Source:
\texttt{scripts/k\_sweep/data\_current/heatmap\_data.csv}.)

\subsubsection{5.3 Failure mode (c) --- (n\_c choose 2) penalty effect
(small-N text
domain)}\label{53-failure-mode-c--n_c-choose-2-penalty-effect-small-n-text-domain}

In LayerForge\textquotesingle s operational range (N=20-40 cross-domain
text corpus):

\begin{itemize}
\tightlist
\item
  edge density is moderate (neither sparse nor dense)
\item
  as the common-reference data of §4.2 shows, the above-limit fraction
  lies in 0.25-0.80, with most cells ≥ 0.5 at N ≥ 20 →
  Newman\textquotesingle s resolution limit is not operationally binding
\item
  when CPM is applied, the \texttt{(n\_c\ choose\ 2)} term relatively
  dominates \texttt{m\_c}, "large clusters are H-loss" becomes the
  dominant force, and CPM selects a \textbf{systematically coarse
  partition} (under-merging that remains even under K-matched
  comparison)
\end{itemize}

This is the root cause of the numerical gaps \textbf{§4.4 (H\_struct
87.5\% vs 19\%)} and \textbf{§4.6 Newman 0.557 vs CPM 0.239 under forced
K=4}. Mode (c) is the component that remains even when K is fixed, and
is the \textbf{structural} primary cause of Newman\textquotesingle s
advantage in this domain.

\subsubsection{5.4 Failure mode (b) --- Calibration bias of
find\_cpm\_resolution}\label{54-failure-mode-b--calibration-bias-of-find_cpm_resolution}

This implementation\textquotesingle s \texttt{find\_cpm\_resolution}
(and the corresponding \texttt{find\_valid\_scale} on the Newman side)
is designed to \textbf{prefer the smallest K within target\_range}.
Specifically, \texttt{cpm\_backend.py::find\_cpm\_resolution}
(L262-263):

\begin{Shaded}
\begin{Highlighting}[]
\ControlFlowTok{if}\NormalTok{ K\_min }\OperatorTok{\textless{}=}\NormalTok{ k }\OperatorTok{\textless{}=}\NormalTok{ K\_max:}
\NormalTok{    best }\OperatorTok{=}\NormalTok{ (mid, labels, h, k)}
\NormalTok{    hi }\OperatorTok{=}\NormalTok{ mid  }\CommentTok{\# prefer smaller γ (smaller K)}
\end{Highlighting}
\end{Shaded}

This design reflects the operational judgment "a coarser partition is
easier to interpret (under the 4±1 constraint)" (if the working-memory
upper bound is an ergonomic upper bound, smaller K is more desirable).
As a side effect, however, \textbf{tracking failure occurs on corpora
where n\_themes does not match target\_range{[}0{]}}.

The phenomenon in §4.6 (20 Newsgroups, n\_themes=4, target\_range=(3,5))
that both methods converge to K=3 is a symptom of this mode. Under
forced K=4 Newman rises to ARI=0.557 (avoiding mode (b) leaves only the
effect of mode (c)), but if we look only at default behavior, Newman
also fails to track n\_themes.

The reason \textbf{Newman achieves 87.5\% tracking in §4.4 (synthetic
corpus, K ranges densely covering around n\_themes)} is that the K-range
design places n\_themes near target\_range{[}0{]}, so the influence of
mode (b) is small. We note in this paper that for external corpora, mode
(b) becomes visible under the default target\_range=(3,5).

\subsubsection{5.5 Failure mode (a) --- Louvain refinement gap
(cross-ref to
§7.3)}\label{55-failure-mode-a--louvain-refinement-gap-cross-ref-to-73}

This implementation\textquotesingle s CPM-Louvain is vanilla Louvain
(greedy single-node moves) and does not include the Leiden refinement
(Traag 2019). On dense small graphs such as Karate Club, the K=2 macro
split is not a local optimum of single-node moves, and even lowering γ
leaves it stuck at K=4 sub-clustering (a known phenomenon in the
literature). §7.3 of this paper makes it explicit that the Karate Club
result (ARI 0.595, K=4) falls under this mode.

For the main results in the small-scale text domain of this
implementation (§4.4, §4.6), the graphs are not dense small graphs and
mode (a) is not dominant. However, the \textbf{phenomenon that K comes
out larger than the ground truth in the external comparison (Karate
Club)} is explained by this mode. It is an algorithm-engineering
constraint, independent of CPM\textquotesingle s penalty structure
itself.

\subsubsection{5.6 Implication --- "resolution-limit-free is universally
better" is
domain-dependent}\label{56-implication--resolution-limit-free-is-universally-better-is-domain-dependent}

CPM\textquotesingle s theoretical advantage manifests in \textbf{domains
where the resolution limit binds}. In LayerForge\textquotesingle s
small-N text domain, as the §4.2 common-reference data shows,
limit-binding is rare in the first place, and CPM\textquotesingle s
"fix" targets \textbf{a problem that is not a problem in this domain}.
At the same time, CPM\textquotesingle s own penalty structure (mode (c))
and calibration design (mode (b)) produce \textbf{different problems},
and Newman consequently exceeds CPM consistently.

This is not a direct refutation of Traag et al. (2011)\textquotesingle s
theoretical result (CPM\textquotesingle s subgraph-invariance). It is
the mechanism observation that \textbf{method selection must be made in
light of the operating-domain characteristics (whether the limit binds,
the typical relation between n\_themes and target\_range, the
dense/sparse degree of the graph)}.

\begin{center}\rule{0.5\linewidth}{0.5pt}\end{center}

\subsubsection{5.7 Empirical method-selection rule
(sketch)}\label{57-empirical-method-selection-rule-sketch}

Given the mechanism analysis in §5.2-§5.5, \textbf{empirical signals}
for choosing "Newman or CPM" at operation time can be constructed:

{\def\LTcaptype{none} % do not increment counter
\begin{longtable}[]{@{}ll@{}}
\toprule\noalign{}
Signal & Recommended threshold \\
\midrule\noalign{}
\endhead
\bottomrule\noalign{}
\endlastfoot
Newman Q at default K (K=4) & Q ≥ 0.3 → Newman \\
Above-limit fraction at default K & ≥ 0.5 → Newman safe \\
Edge density at θ & \textgreater{} 0.5 → consider CPM \\
Corpus size N & \textgreater{} 1000 → consider CPM \\
Cross-method ARI at probe K & \textless{} 0.3 → flag uncertainty \\
\end{longtable}
}

Pseudocode for the decision rule is detailed in the method-selection
rule section of the companion document (GitHub repository:
\url{https://github.com/ChaiCroquis/LayerForge/blob/main/docs/08_empirical_findings.md}).
\textbf{We have not implemented it as code}: in
LayerForge\textquotesingle s operational range (N=20-40, moderate edge
density), the rule outputs "newman" in \textasciitilde100\% of cases, so
the implementation cost outweighs the information value. The design is
to be used as a foundation when extending to other domains (N=1000+,
dense graphs).

\begin{center}\rule{0.5\linewidth}{0.5pt}\end{center}

\subsection{§6. Application examples}\label{6-application-examples}

LayerForge\textquotesingle s Modes A/B/C have application areas
corresponding to each of the three costs of §1 (cognitive load /
bandwidth / AI cost):

\subsubsection{6.1 Mode A --- Reducing the reader\textquotesingle s
cognitive load}\label{61-mode-a--reducing-the-readers-cognitive-load}

Decompose a corpus into 4±1 layers, and expand only the "layer(s) of
interest" from the UI / CLI. The reader can judge the content from the
representative expression of the layer of interest, without reading the
full corpus.

\subsubsection{6.2 Mode B --- Managing concern transitions in decision
information}\label{62-mode-b--managing-concern-transitions-in-decision-information}

Layer-organize decision information under consideration, and attach
\texttt{open}/\texttt{close}/\texttt{settle} states. By separating items
under decision (open) from items already decided (settled), the reader
can focus only on the \textbf{currently activated concern}.

\subsubsection{6.3 Mode C --- Multi-agent bandwidth
reduction}\label{63-mode-c--multi-agent-bandwidth-reduction}

Compress AI agent A\textquotesingle s verbose output (e.g., a 15K-token
comprehensive answer) into the layer subset relevant to agent
B\textquotesingle s query (decision-less subset guarantee --- i.e., it
functions only as an information subset, with no LLM-side
summarization). In this study, \texttt{scripts/multiagent\_demo/}
achieves a 63\% context reduction (B-axis null result: the downstream
agent\textquotesingle s output quality is preserved even after
compression).

\textbf{Relation to §4.5\textquotesingle s 1/K=10\% (different
measurements)}: The "K=10 with 1/K=10\% (10× reduction)" shown in §4.5
is the \textbf{theoretical minimum read-fraction on the
partition-consistency metric} (the upper bound saying that, assuming
self-routing of a passage-as-query succeeds, reading only the one layer
it belongs to is sufficient); the unit of measurement is "number of
layers read / total layers". The 63\% reduction of Mode C in this
section is a \textbf{token-count-based empirical value}, comparing agent
A\textquotesingle s full output and the subset that agent B actually
reads, in token units, in the demo script. The two are \textbf{different
measurements} (theoretical minimum vs. production-empirical) and cannot
be directly compared. The former is the floor on a fine-grained
partition at K=10; the latter is the measured reduction in a realistic
Mode C workflow.

\begin{center}\rule{0.5\linewidth}{0.5pt}\end{center}

\subsection{§7. Limitations}\label{7-limitations}

\subsubsection{7.1 Scope of validation}\label{71-scope-of-validation}

\begin{itemize}
\tightlist
\item
  \textbf{The operational target is the N=20-40 small-N text domain.} As
  an external benchmark, we scale up to N=100 on 20 Newsgroups (§4.6),
  and the conclusion goes in the same direction (Newman advantage).
  \textbf{The ARI table in §4.3 is N=12-40 and §4.6 is N=100, so the
  intermediate N range (40-100) is not directly swept in this paper},
  but since Newman\textquotesingle s advantage is observed in the same
  direction at both endpoints, we infer the behavior in the interpolated
  range to be in the same direction (confirmation is future work).
  \textbf{N \textgreater{} 1000 and dense graphs are unverified.}
\item
  \textbf{The external benchmark is only one corpus (20NG)} ---
  reproduction on other public corpora is future work.
\item
  \textbf{CPM-Louvain has no Leiden refinement} --- corresponds to the
  known Karate Club K=4 sub-clustering phenomenon in the literature;
  specifically positioned in §7.3.
\end{itemize}

\subsubsection{7.2 B-axis null results}\label{72-b-axis-null-results}

\textbf{Improvement in the LLM\textquotesingle s own behavior} when
LayerForge is integrated into an AI workflow is outside the scope of
this study; the three trials we ran (hallucination benchmark,
multi-agent drift verification, context filter ablation) all gave
\textbf{null results}. This means the claim "LayerForge makes AI
smarter" is not defensible. The contribution of this study lies in
\textbf{the pipeline itself that narrows AI output by concern and hands
it to the reader / downstream agent}; we do not claim improvement in AI
behavior.

\subsubsection{7.3 Correctness gate for the
self-implementation}\label{73-correctness-gate-for-the-self-implementation}

A head-to-head numerical-accuracy comparison against the reference
implementation of CPM-Louvain (GPL \texttt{leidenalg}) was not performed
due to MIT-license-compatibility issues. As a substitute, we put in
place the following \textbf{three-stage indirect gate}:

\begin{enumerate}
\def\labelenumi{\arabic{enumi}.}
\tightlist
\item
  \textbf{Synthetic 3-block test}
  (\texttt{tests/axioms/test\_cpm\_backend.py::test\_cpm\_partition\_separates\_three\_blocks}):
  K=3 is perfectly recovered on 3 blocks (intra=0.9, inter=0.05).
\item
  \textbf{Karate Club ARI = 0.595} vs the empirical 2-community ground
  truth (chance baseline = 0): on the public Zachary (1977) graph,
  because of over-splitting due to the absence of the Leiden refinement
  (the K=4 sub-clustering phenomenon mentioned in §7.1), directly
  inspecting the partition gives K=4 clusters. From the cluster
  structure we can verify that this is a \textbf{sub-refinement} of the
  ground-truth 2-community split, and ARI=0.595 is an intermediate value
  indicating \textbf{the agreement between this sub-refinement and the
  ground truth} (perfect match would give ARI=1.0, random would give 0).
  This suggests that adding the Leiden refinement leaves room to reach
  the K=2 macro split.
\item
  \textbf{Determinism cross-check}: both Newman and CPM are reproducible
  under the same seed.
\end{enumerate}

This is not a strong proof of "CPM is implemented correctly" (the
local-optimum limitation in the absence of the Leiden refinement is
mentioned in §5.5 / §7.1). It is positioned as \textbf{evidence that the
minimum correctness requirement is met as a prerequisite for the
empirical comparison against Newman}.

\subsubsection{7.4 No universality claim for
4±1}\label{74-no-universality-claim-for-41}

We adopt Cowan (2010)\textquotesingle s working-memory capacity as a
design target of this study, but we do not claim that "4±1 universally
emerges on any corpus". In our measurements, K\_optimal fluctuates in
2-11 depending on the corpus (§4.1); we keep its positioning as an
ergonomic upper bound.

\subsubsection{7.5 Ablations declared out of
scope}\label{75-ablations-declared-out-of-scope}

This paper restricts its scope to validating the
\textbf{community-detection method choice (Newman vs CPM)}. The
following ablations on other choices within the combination are not
performed:

\begin{itemize}
\tightlist
\item
  The effect of \textbf{SCA distillation on/off} on partition / output
  quality
\item
  The effect of \textbf{HERCULES hierarchy depth = 1 (this
  paper\textquotesingle s default) vs depth \textgreater{} 1}
\item
  The difference in cognitive-load reduction between \textbf{the 4±1
  constraint vs. no constraint (arbitrary K\_optimal)}
\item
  Output-drift measurement under \textbf{LLM involvement on/off at
  Boundary 1/2}
\end{itemize}

These would be needed to validate the LayerForge combination as a whole,
more broadly than this paper does. They lie outside the claim range of
this paper (validity of the most important choice inside the
combination, the community-detection method) and are left as future
work.

\begin{center}\rule{0.5\linewidth}{0.5pt}\end{center}

\subsection{§8. Conclusions}\label{8-conclusions}

For the problem stated in §1 --- \textbf{converging AI output by concern
to reduce cognitive load / bandwidth / cost} --- we showed in §3 that
the five-component combination of §2 is a \textbf{design that satisfies
the task requirements R1-R5}, and in §4 we empirically tested the most
important choice within the combination (the community-detection
method).

Main observations (within this domain and within the verification range
of this paper):

\begin{enumerate}
\def\labelenumi{\arabic{enumi}.}
\tightlist
\item
  \textbf{The above-limit fraction (the Fortunato-Barthélemy ratio)
  functions as a supplementary K-selection signal that does not depend
  on the community-detection method} (§4.2).
\item
  \textbf{Within the test range (16 settings + 20 Newsgroups at N=100),
  the Newman backend exceeded the CPM backend by a large effect size}
  (H\_struct 87.5\% vs 19\%; 20NG ARI 0.430 vs 0.239 at default K=3;
  Newman reaches 0.557 when K=4 is forced; §4.4 / §4.6). \textbf{Because
  of the small sample, we do not perform formal statistical significance
  tests and report qualitatively by effect size.}
\item
  \textbf{CPM\textquotesingle s underperformance is mechanistically
  explainable}: the \texttt{(n\_c\ choose\ 2)} quadratic penalty induces
  under-merging. The proposition "resolution-limit-free is universally
  better" does not hold domain-independently (§5).
\item
  \textbf{On the self-routing metric (a partition-consistency test that
  uses each passage as its own query), the theoretical minimum
  read-fraction of 1/K=10\% at K=10 holds for both methods} (§4.5). This
  is not a measure of production query-routing accuracy; it reports that
  the \textbf{necessary condition for partition consistency} is
  satisfied. Production paraphrased-query routing accuracy, and the
  relation to the 63\% empirical token reduction reported for Mode C in
  §6.3, are stated in §4.5 / §6.3; production accuracy verification is
  future work (§7.2 / §7.5).
\item
  \textbf{A method-selection rule can be constructed from empirical
  signals} (§5.7; details in the companion document on the GitHub
  repository:
  \url{https://github.com/ChaiCroquis/LayerForge/blob/main/docs/08_empirical_findings.md}).
  In this domain the rule outputs "newman" almost always, so we do not
  implement it at runtime and preserve it only in the docs.
\end{enumerate}

LayerForge is released under the MIT license, can be operated as a
Claude Code skill, and bundles reproducibility scripts and CSV / JSON.
The numbers measured in this paper are pinned to the publication-time
commit hash (see \texttt{git\ log}).

\begin{center}\rule{0.5\linewidth}{0.5pt}\end{center}

\subsection{References}\label{references}

\subsubsection{Foundational}\label{foundational}

\begin{itemize}
\tightlist
\item
  Cowan, N. (2010). The magical mystery four: How is working memory
  capacity limited, and why? \emph{Current Directions in Psychological
  Science}, 19(1), 51-57.
\item
  Zachary, W. W. (1977). An information flow model for conflict and
  fission in small groups. \emph{Journal of Anthropological Research},
  33(4), 452-473.
\end{itemize}

\subsubsection{Community detection theory and
benchmarks}\label{community-detection-theory-and-benchmarks}

\begin{itemize}
\tightlist
\item
  Newman, M. E. J. (2006). Modularity and community structure in
  networks. \emph{PNAS}, 103(23), 8577-8582.
\item
  Fortunato, S., \& Barthélemy, M. (2007). Resolution limit in community
  detection. \emph{PNAS}, 104(1), 36-41.
\item
  Lancichinetti, A., \& Fortunato, S. (2009). Benchmarks for testing
  community detection algorithms on directed and weighted graphs with
  overlapping communities. \emph{Physical Review E}, 80, 016118.
\item
  Good, B. H., de Montjoye, Y.-A., \& Clauset, A. (2010). Performance of
  modularity maximization in practical contexts. \emph{Physical Review
  E}, 81, 046106.
\item
  Traag, V. A., Van Dooren, P., \& Nesterov, Y. (2011). Narrow scope for
  resolution-limit-free community detection. \emph{Physical Review E},
  84, 016114.
\item
  Aldecoa, R., \& Marín, I. (2013). Exploring the limits of community
  detection strategies in complex networks. \emph{Scientific Reports},
  3, 2216.
\item
  Traag, V. A., Waltman, L., \& van Eck, N. J. (2019). From Louvain to
  Leiden: guaranteeing well-connected communities. \emph{Scientific
  Reports}, 9, 5233.
\item
  Felipe, L. L., Avrachenkov, K., \& Menasche, D. S. (2025). From Leiden
  to Pleasure Island: The Constant Potts Model for community detection
  as a hedonic game. arXiv:2509.03834.
\end{itemize}

\subsubsection{Clustering-based topic modeling and graph-based text
clustering}\label{clustering-based-topic-modeling-and-graph-based-text-clustering}

\begin{itemize}
\tightlist
\item
  Angelov, D. (2020). Top2Vec: Distributed representations of topics.
  arXiv:2008.09470.
\item
  Rao, R. N., \& Chakraborty, M. (2021). Vec2GC -- A graph based
  clustering method for text representations. arXiv:2104.09439.
\item
  Grootendorst, M. (2022). BERTopic: Neural topic modeling with a
  class-based TF-IDF procedure. arXiv:2203.05794.
\item
  Zhang, Z., Fang, M., Chen, L., \& Namazi-Rad, M.-R. (2022). Is neural
  topic modelling better than clustering? An empirical study on
  clustering with contextual embeddings for topics. \emph{Proceedings of
  NAACL}. (CETopic)
\item
  Zhang, L., Liu, J., \& Yan, Q. (2023). Graph2Topic: An opensource
  topic modeling framework based on sentence embedding and community
  detection. arXiv:2304.06653. (G2T)
\item
  Eichin, F., Schuster, M., Groh, G., \& Hedderich, M. A. (2024).
  Semantic Component Analysis: Discovering patterns in short texts
  beyond topics. arXiv:2410.21054.
\item
  Janssens, W., Bogaert, M., \& Van den Poel, D. (2025). LLM-assisted
  topic reduction for BERTopic on social media data. arXiv:2509.19365.
\end{itemize}

\subsubsection{Hierarchical clustering, RAG and
GraphRAG}\label{hierarchical-clustering-rag-and-graphrag}

\begin{itemize}
\tightlist
\item
  Sarthi, P., Abdullah, S., Tuli, A., Khanna, S., Goldie, A., \&
  Manning, C. D. (2024). RAPTOR: Recursive Abstractive Processing for
  Tree-Organized Retrieval. \emph{ICLR}. arXiv:2401.18059.
\item
  Edge, D., Trinh, H., Cheng, N., Bradley, J., Chao, A., Mody, A.,
  Truitt, S., \& Larson, J. (2024). From local to global: A Graph RAG
  approach to query-focused summarization. arXiv:2404.16130.
\item
  Xiao, Y., Dong, J., Zhou, C., Dong, S., Zhang, Q.-w., Yin, D., Sun,
  X., \& Huang, X. (2025). GraphRAG-Bench: Challenging domain-specific
  reasoning for evaluating graph retrieval-augmented generation.
  arXiv:2506.02404.
\item
  Han, H., Ma, L., Wang, Y., Shomer, H., Lei, Y., Qi, Z., Guo, K., Hua,
  Z., Long, B., Liu, H., Aggarwal, C. C., \& Tang, J. (2026). RAG vs.
  GraphRAG: A systematic evaluation and key insights. arXiv:2502.11371.
\item
  Hossain, J., \& Sarıyüce, A. E. (2026). Core-based hierarchies for
  efficient GraphRAG. arXiv:2603.05207.
\end{itemize}

\subsubsection{Context and memory
compression}\label{context-and-memory-compression}

\begin{itemize}
\tightlist
\item
  Bohdal, O., Saha, P., Michieli, U., Ozay, M., \& Ceritli, T. (2026).
  Clustering-driven memory compression for on-device large language
  models. arXiv:2601.17443.
\item
  Zhou, Y., Lei, Y., Si, S., Sun, Q., Wang, W., Wu, Y., Wen, H., Chen,
  G., Qi, F., \& Sun, M. (2025). From context to EDUs: Faithful and
  structured context compression via Elementary Discourse Unit
  decomposition. arXiv:2512.14244.
\end{itemize}

\subsubsection{Reference
implementations}\label{reference-implementations}

\begin{itemize}
\tightlist
\item
  \texttt{mainlp/semantic\_components} (MIT, 2024) --- SCA reference
  implementation by F. Eichin
\item
  \texttt{bandeerun/pyhercules} (MIT, 2025) --- HERCULES Python
  implementation (no peer-reviewed publication at the time of this work)
\end{itemize}

\subsubsection{Project artifact}\label{project-artifact}

\begin{itemize}
\tightlist
\item
  LayerForge: \url{https://github.com/ChaiCroquis/LayerForge} (MIT
  licensed)
\item
  Commit hash for measurements in this paper: see \texttt{git\ log} at
  submission time
\item
  All CSVs, JSONs, PNGs in \texttt{scripts/k\_sweep/} are reproducible
  from the corresponding \texttt{*.py} scripts.
\end{itemize}

\begin{center}\rule{0.5\linewidth}{0.5pt}\end{center}

\subsection{Appendix A --- Related work summary
table}\label{appendix-a--related-work-summary-table}

A reference table for the literature detailed in §1.2 (Related work and
positioning). This appendix supplements §1.2; for the detailed
positioning of each work, see §1.2.

{\def\LTcaptype{none} % do not increment counter
\begin{longtable}[]{@{}llll@{}}
\toprule\noalign{}
Group (§1.2 ID) & Main works & Common ground & Differentiation from
LayerForge \\
\midrule\noalign{}
\endhead
\bottomrule\noalign{}
\endlastfoot
Clustering-based topic modeling and graph-based text clustering (§1.2.1)
& BERTopic (Grootendorst 2022), Top2Vec (Angelov 2020), CETopic (Zhang
et al. NAACL 2022), Vec2GC (Rao \& Chakraborty 2021), \textbf{G2T}
(Zhang et al. 2023), SCA (Eichin et al. 2024), LLM-Assisted Topic
Reduction for BERTopic (Janssens et al. 2025) & document/passage
embedding + clustering + extracting topic / cluster structure & G2T is
the closest substrate; density-based works do not use CD. LayerForge
combines SCA distillation + HERCULES skeleton + 4±1 + Mode A/B/C, and
adds a systematic Newman-vs-CPM comparison (§4-§5). \\
Hierarchical RAG and GraphRAG (§1.2.2) & RAPTOR (Sarthi et al. 2024),
GraphRAG (Edge et al. 2024), HippoRAG, \textbf{Core-based GraphRAG}
(Hossain \& Sarıyüce 2026), RAG vs. GraphRAG eval (Han et al. 2026),
GraphRAG-Bench (Xiao et al. 2025), HERCULES
(\texttt{bandeerun/pyhercules} 2025) & hierarchical structure +
retrieval/summarization; many assume LLM integration & LayerForge
decouples LLM from the analysis core and prioritizes reproducibility
(R5). Shares "Leiden\textquotesingle s reproducibility problem"
recognition with Core-based GraphRAG, but pursues the
failure-mode-documentation direction rather than replacement
(k-core). \\
Context and memory compression (§1.2.3) & \textbf{Clustering-driven
Memory Compression} (Bohdal et al. 2026, Samsung), EDU Context
Compressor (Zhou et al. 2025) & clustering / structure-based context
compression & Bohdal differentiates on target (memory tokens vs. corpus
passages) / goal (personalization vs. concern-wise convergence) / method
(averaging-merge vs. community detection) / cognitive constraint (none
vs. Cowan 4±1) --- detailed in §1.2.5 \#2. EDU is discourse-structural
and orthogonal to this work. \\
Community detection method comparison (§1.2.4) & Lancichinetti \&
Fortunato (2009), Aldecoa \& Marín (2013), Traag et al. (2011, 2019),
\textbf{From Leiden to Pleasure Island} (Felipe et al. 2025) &
systematic comparison / theoretical analysis of community-detection
algorithms & Prior works center on synthetic / random / closed
benchmarks. LayerForge provides empirical evaluation on small-N text
similarity graphs (N=20-40). Felipe et al. (Pleasure Island) complements
§5 of this work via game-theoretic analysis of CPM. \\
\end{longtable}
}

LayerForge\textquotesingle s contribution narrows to \textbf{(i) an
implementation report of combining known techniques (G2T-like substrate
+ SCA + HERCULES + 4±1 + Mode A/B/C), and (ii) a systematic
Newman-vs-CPM comparison in the small-N text domain + a
three-failure-mode classification}. We do not claim \textbf{the
community-detection method comparison itself} as novelty; the paper is
positioned as \textbf{a report that contributes empirical data in this
domain to the literature space}. The concurrent framing overlap with
Bohdal et al. (Samsung 2026) does not negate the novelty of this work;
we record it as the fact that two studies independently identified the
same problem space.

\subsection{Appendix B --- Sensitivity analyses
summary}\label{appendix-b--sensitivity-analyses-summary}

A summary of the main outcomes of the sensitivity analyses detailed from
§2 onward of the companion document on the GitHub repository
(\url{https://github.com/ChaiCroquis/LayerForge/blob/main/docs/08_empirical_findings.md}):

\begin{itemize}
\tightlist
\item
  Embedding model (MiniLM vs mpnet): mpnet gives a cleaner Q peak; both
  methods go in the same direction.
\item
  Random seed (42 / 123): K\_optimal is stable; Newman\textquotesingle s
  Q peak K is seed-invariant.
\item
  Edge floor: above-limit fraction is monotonic in the same way under 0
  (default) and median similarity (the CPM reference).
\item
  N (corpus size): K=n\_themes tracking is stable for Newman from N≥20;
  CPM is weak and N-insensitive.
\end{itemize}

\begin{center}\rule{0.5\linewidth}{0.5pt}\end{center}

\emph{Document ends.}

\end{document}
